% GENERATED FILE. Edit canonical lessons, then run npm run generate:books.
% canonical-source-hash: fnv1a64:fe11a7bbcc64a3e9
% canonical-lessons: TA-C489-uranku, TA-C489-kurattai-vitu, TA-C489-pural, TA-C489-vili, TA-C489-viyar

\chapter{A Restless Night}
\label{ch:ta-a-restless-night}
\hlchaptermodality{489}

\begin{chapteropening}
\textbf{By the end of this chapter:} \emph{I can talk about slumber, snoring, tossing and turning, waking and sweating.}

Complete the last lesson of chapter 489: I can talk about slumber, snoring, tossing and turning, waking and sweating.
\end{chapteropening}

\section[\texorpdfstring{\ta{உறங்கு} (uṟaṅku)}{உறங்கு (uṟaṅku)}]{\ta{உறங்கு} (uṟaṅku) — to slumber}

\label{lesson:TA-C489-uranku}



\begin{quote}
Before the new one: say the Tamil for a hunch, then the Tamil for starvation.
\end{quote}

\subsection*{You'll want to know: \ta{உறங்கு}}
\textbf{\ta{உறங்கு}} — \emph{uṟaṅku} — \textquotedblleft{}to slumber\textquotedblright{}.

\begin{quote}
\ta{குழந்தை} \ta{அமைதியாக} \ta{உறங்கியது}.
\end{quote}

\begin{quote}
\emph{kuḻantai amaitiyāka uṟaṅkiyatu.}
\end{quote}

\begin{quote}
The baby slept peacefully.
\end{quote}

Past (I …): \ta{உறங்கினேன்}, \emph{uṟaṅkiṉēṉ}.

A softer, more written word than \ta{தூங்கு}.

\begin{cousinweb}
\textbf{\ta{உறங்கு}} (\emph{uṟaṅku}) is native Dravidian (DEDR 707); the entry also gives \ta{உறக்கு} \emph{uṟakku}, \textquotedblleft{}to put to sleep\textquotedblright{}.
\end{cousinweb}

\begin{grammarlens}[title={What you've built}]
One more word for describing what you lived and felt.
\end{grammarlens}

\subsection*{Your turn}
\begin{itemize}
\raggedright
  \item \emph{Say it:} \emph{uṟaṅku}
  \item \emph{Say it:} \emph{uṟaṅku}, once more
  \item \emph{Say it:} \emph{kuḻantai amaitiyāka uṟaṅkiyatu.} — The baby slept peacefully.
  \item \emph{Recall:} say the Tamil for a hunch, then the Tamil for starvation, then say \emph{uṟaṅku} again
\end{itemize}

\begin{tcolorbox}[breakable,colback=teal!4,colframe=teal!35!black,title={Before you move on}]
What does \ta{உறங்கு} mean? (\textquotedblleft{}To slumber\textquotedblright{}.) Say it once more.
\end{tcolorbox}
\section[\texorpdfstring{\ta{குறட்டை} \ta{விடு} (kuṟaṭṭai viṭu)}{குறட்டை விடு (kuṟaṭṭai viṭu)}]{\ta{குறட்டை} \ta{விடு} (kuṟaṭṭai viṭu) — to snore}

\label{lesson:TA-C489-kurattai-vitu}



\begin{quote}
Before the new one: say the Tamil for starvation, then the Tamil for to slumber.
\end{quote}

\subsection*{You'll want to know: \ta{குறட்டை} \ta{விடு}}
\textbf{\ta{குறட்டை} \ta{விடு}} — \emph{kuṟaṭṭai viṭu} — \textquotedblleft{}to snore\textquotedblright{}.

\begin{quote}
\ta{அப்பா} \ta{இரவு} \ta{முழுவதும்} \ta{குறட்டை} \ta{விட்டார்}.
\end{quote}

\begin{quote}
\emph{appā iravu muḻuvatum kuṟaṭṭai viṭṭār.}
\end{quote}

\begin{quote}
Father snored all night.
\end{quote}

Past (I …): \ta{குறட்டை} \ta{விட்டேன்}, \emph{kuṟaṭṭai viṭṭēṉ}.

\ta{குறட்டை} \textquotedblleft{}snoring\textquotedblright{} + \ta{விடு} \textquotedblleft{}to let out\textquotedblright{}.

\begin{cousinweb}
\textbf{\ta{குறட்டை} \ta{விடு}} (\emph{kuṟaṭṭai viṭu}) is native: \ta{குறட்டை} \emph{kuṟaṭṭai}, \textquotedblleft{}snoring\textquotedblright{} (DEDR 1852), with \ta{விடு} \textquotedblleft{}to let out\textquotedblright{}.
\end{cousinweb}

\begin{grammarlens}[title={What you've built}]
One more word for describing what you lived and felt.
\end{grammarlens}

\subsection*{Your turn}
\begin{itemize}
\raggedright
  \item \emph{Say it:} \emph{kuṟaṭṭai viṭu}
  \item \emph{Say it:} \emph{kuṟaṭṭai viṭu}, once more
  \item \emph{Say it:} \emph{appā iravu muḻuvatum kuṟaṭṭai viṭṭār.} — Father snored all night.
  \item \emph{Recall:} say the Tamil for starvation, then the Tamil for to slumber, then say \emph{kuṟaṭṭai viṭu} again
\end{itemize}

\begin{tcolorbox}[breakable,colback=teal!4,colframe=teal!35!black,title={Before you move on}]
What does \ta{குறட்டை} \ta{விடு} mean? (\textquotedblleft{}To snore\textquotedblright{}.) Say it once more.
\end{tcolorbox}
\section[\texorpdfstring{\ta{புரள்} (puraḷ)}{புரள் (puraḷ)}]{\ta{புரள்} (puraḷ) — to toss and turn; to roll}

\label{lesson:TA-C489-pural}



\begin{quote}
Before the new one: say the Tamil for to slumber, then the Tamil for to snore.
\end{quote}

\subsection*{You'll want to know: \ta{புரள்}}
\textbf{\ta{புரள்}} — \emph{puraḷ} — \textquotedblleft{}to toss and turn; to roll\textquotedblright{}.

\begin{quote}
\ta{தூக்கம்} \ta{வராமல்} \ta{நான்} \ta{இரவு} \ta{முழுவதும்} \ta{புரண்டேன்}.
\end{quote}

\begin{quote}
\emph{tūkkam varāmal nāṉ iravu muḻuvatum puraṇṭēṉ.}
\end{quote}

\begin{quote}
Unable to sleep, I tossed and turned all night.
\end{quote}

Past (I …): \ta{புரண்டேன்}, \emph{puraṇṭēṉ}.

Past \ta{புரண்டேன்}, with \ta{ள்} + \ta{ந்} written \ta{ண்ட}.

\begin{cousinweb}
\textbf{\ta{புரள்}} (\emph{puraḷ}) is native Dravidian (DEDR 4285), of rolling waves and rolling bodies.
\end{cousinweb}

\begin{grammarlens}[title={What you've built}]
One more word for describing what you lived and felt.
\end{grammarlens}

\subsection*{Your turn}
\begin{itemize}
\raggedright
  \item \emph{Say it:} \emph{puraḷ}
  \item \emph{Say it:} \emph{puraḷ}, once more
  \item \emph{Say it:} \emph{tūkkam varāmal nāṉ iravu muḻuvatum puraṇṭēṉ.} — Unable to sleep, I tossed and turned all night.
  \item \emph{Recall:} say the Tamil for to slumber, then the Tamil for to snore, then say \emph{puraḷ} again
\end{itemize}

\begin{tcolorbox}[breakable,colback=teal!4,colframe=teal!35!black,title={Before you move on}]
What does \ta{புரள்} mean? (\textquotedblleft{}To toss and turn; to roll\textquotedblright{}.) Say it once more.
\end{tcolorbox}
\section[\texorpdfstring{\ta{விழி} (viḻi)}{விழி (viḻi)}]{\ta{விழி} (viḻi) — to wake up; to open one’s eyes}

\label{lesson:TA-C489-vili}



\begin{quote}
Before the new one: say the Tamil for to snore, then the Tamil for to toss and turn.
\end{quote}

\subsection*{You'll want to know: \ta{விழி}}
\textbf{\ta{விழி}} — \emph{viḻi} — \textquotedblleft{}to wake up; to open one’s eyes\textquotedblright{}.

\begin{quote}
\ta{நள்ளிரவில்} \ta{சத்தம்} \ta{கேட்டு} \ta{நான்} \ta{விழித்தேன்}.
\end{quote}

\begin{quote}
\emph{naḷḷiravil cattam kēṭṭu nāṉ viḻittēṉ.}
\end{quote}

\begin{quote}
I woke up at midnight on hearing a noise.
\end{quote}

Past (I …): \ta{விழித்தேன்}, \emph{viḻittēṉ}.

As a noun, the eye: \ta{விழிகள்}.

\begin{cousinweb}
\textbf{\ta{விழி}} (\emph{viḻi}) is native Dravidian (DEDR 5429), both the eye and the opening of it.
\end{cousinweb}

\begin{grammarlens}[title={What you've built}]
One more word for describing what you lived and felt.
\end{grammarlens}

\subsection*{Your turn}
\begin{itemize}
\raggedright
  \item \emph{Say it:} \emph{viḻi}
  \item \emph{Say it:} \emph{viḻi}, once more
  \item \emph{Say it:} \emph{naḷḷiravil cattam kēṭṭu nāṉ viḻittēṉ.} — I woke up at midnight on hearing a noise.
  \item \emph{Recall:} say the Tamil for to snore, then the Tamil for to toss and turn, then say \emph{viḻi} again
\end{itemize}

\begin{tcolorbox}[breakable,colback=teal!4,colframe=teal!35!black,title={Before you move on}]
What does \ta{விழி} mean? (\textquotedblleft{}To wake up; to open one’s eyes\textquotedblright{}.) Say it once more.
\end{tcolorbox}
\section[\texorpdfstring{\ta{வியர்} (viyar)}{வியர் (viyar)}]{\ta{வியர்} (viyar) — to perspire; to sweat}

\label{lesson:TA-C489-viyar}



\begin{quote}
Before the new one: say the Tamil for to toss and turn, then the Tamil for to wake up.
\end{quote}

\subsection*{You'll want to know: \ta{வியர்}}
\textbf{\ta{வியர்}} — \emph{viyar} — \textquotedblleft{}to perspire; to sweat\textquotedblright{}.

\begin{quote}
\ta{வெயிலில்} \ta{நடந்ததால்} \ta{எனக்கு} \ta{வியர்த்தது}.
\end{quote}

\begin{quote}
\emph{veyilil naṭantatāl eṉakku viyarttatu.}
\end{quote}

\begin{quote}
I sweated because I walked in the sun.
\end{quote}

Past (it …): \ta{வியர்த்தது}, \emph{viyarttatu}.

Dative experiencer: \ta{எனக்கு} \ta{வியர்த்தது}. \ta{வியர்வை}, sweat.

\begin{cousinweb}
\textbf{\ta{வியர்}} (\emph{viyar}) is native Dravidian (DEDR 5486), which also lists the form \ta{வெயர்} \emph{veyar}.
\end{cousinweb}

\begin{grammarlens}[title={What you've built}]
One more word for describing what you lived and felt.
\end{grammarlens}

\subsection*{Your turn}
\begin{itemize}
\raggedright
  \item \emph{Say it:} \emph{viyar}
  \item \emph{Say it:} \emph{viyar}, once more
  \item \emph{Say it:} \emph{veyilil naṭantatāl eṉakku viyarttatu.} — I sweated because I walked in the sun.
  \item \emph{Recall:} say the Tamil for to toss and turn, then the Tamil for to wake up, then say \emph{viyar} again
\end{itemize}

\begin{tcolorbox}[breakable,colback=teal!4,colframe=teal!35!black,title={Before you move on}]
What does \ta{வியர்} mean? (\textquotedblleft{}To perspire; to sweat\textquotedblright{}.) Say it once more.
\end{tcolorbox}
